\documentclass[11pt]{article}

\usepackage{amsmath,amssymb,graphicx}
\usepackage{geometry}
\geometry{margin=1in}

\title{Hedging Personal Income Risk Using Prediction Markets}
\author{Rahul Jayachandran, Ethan Lin, Rohan Malhotra}
\date{\today}

\begin{document}

\maketitle

\begin{abstract}
Income risk is one of the largest sources of financial uncertainty for individuals. 
This paper proposes a framework for hedging personal job-loss risk using event 
contracts from prediction markets. Employment is modeled as a stochastic process 
with a macro-sensitive hazard rate governing job loss events. We construct a hedge 
portfolio using contracts whose payoffs are positively correlated with unemployment 
or recession risk. Using Monte Carlo simulations, we evaluate how such hedges reduce 
the left-tail risk of personal income distributions while preserving expected income.
\end{abstract}

\section{Introduction}

Household income risk is typically addressed through savings or traditional 
insurance products. However, these approaches provide limited protection against 
macroeconomic shocks that simultaneously increase unemployment risk and reduce 
future earning opportunities. 

Prediction markets such as Kalshi offer tradable contracts whose payouts are 
tied to macroeconomic events including unemployment levels, recessions, and 
interest rate changes. These instruments provide a potential mechanism for 
hedging personal income risk.

This paper models employment risk as a stochastic process and evaluates how 
prediction market contracts can reduce tail risk in individual income streams.

\section{Income Process}

Let $E_t$ denote the employment state:

\[
E_t =
\begin{cases}
1 & \text{if employed} \\
0 & \text{if unemployed}
\end{cases}
\]

Let $S$ represent salary per unit time. The income accumulation process is

\[
dW_t = S E_t dt
\]

where $W_t$ is cumulative wealth from labor income.


\section{Job Loss Model}

Job loss events are modeled as a stochastic jump process. Let $N_t$ be a Poisson 
process with intensity $\lambda_t$ representing the job-loss hazard rate.

\[
N_t \sim Poisson(\lambda_t)
\]

Employment evolves according to

\[
dE_t = -E_t dN_t
\]

meaning employment transitions from $1$ to $0$ when a job-loss event occurs.

The hazard rate depends on macroeconomic variables:

\[
\lambda_t = \alpha + \beta_1 U_t + \beta_2 r_t + \beta_3 I_t
\]

where

\begin{itemize}
\item $U_t$ = unemployment rate
\item $r_t$ = interest rate
\item $I_t$ = industry-specific risk indicators
\end{itemize}

\section{Income Jump-Diffusion Model}

To capture salary variability and sudden job loss shocks, income is modeled as a 
jump-diffusion process:

\[
dW_t = S dt + \sigma dB_t - J dN_t
\]

where

\begin{itemize}
\item $B_t$ is Brownian motion representing income variability
\item $\sigma$ is salary volatility
\item $J$ is the income loss associated with job loss
\item $N_t$ is the job-loss jump process
\end{itemize}

\section{Prediction Market Hedge}

Let $H_t$ denote the value of a hedge constructed using prediction market 
contracts tied to macroeconomic events.

The combined wealth process becomes

\[
dW_t = S E_t dt + dH_t
\]

The hedge is chosen such that

\[
Corr(\text{JobLoss}, H_t) > 0
\]

ensuring the hedge pays out when unemployment risk increases.

\section{Optimal Hedge Ratio}

We determine the optimal hedge ratio using the minimum variance hedge:

\[
h^* =
\frac{Cov(L, H)}{Var(H)}
\]

where

\begin{itemize}
\item $L$ is income loss due to unemployment
\item $H$ is hedge payoff
\end{itemize}

The optimal hedge reduces variance while minimizing hedge cost.

\section{Model Inputs and Data}

The model requires three categories of inputs.

\subsection{Personal Risk Factors}

These estimate the individual's job-loss probability. Inputs include:
industry, company size, job level (entry, manager, exec), tenure, salary,
Monte Carlo path count, and horizon years. These affect the baseline hazard rate $\lambda(t)$.

\subsection{Macroeconomic Variables}

These drive layoffs and hiring cycles. The state vector is
\[
X_t = (U_t, r_t, \text{GDP}_t, \text{Layoffs}_t)
\]
where $U_t$ is unemployment rate, $r_t$ is the interest rate, $\text{GDP}_t$ 
is growth, and $\text{Layoffs}_t$ is industry layoff activity. Data sources:
\begin{itemize}
\item \textbf{FRED} (Federal Reserve): unemployment (UNRATE), initial jobless 
claims (ICSA), federal funds rate (FEDFUNDS), GDP growth
\item \textbf{BLS JOLTS}: layoffs (JTSLDL), hires (JTSHIR)
\end{itemize}

\subsection{Market Expectations (Forward-Looking Risk)}

Sources: Kalshi API, Polymarket, Fed futures (CME FedWatch), yield curve.
Markets contain implied probabilities: $P(\text{recession})$, 
$P(\text{unemployment} > x)$. FRED provides the 10Y-2Y Treasury spread 
(T10Y2Y) as a recession indicator; inversion historically precedes recessions.

\subsection{Hedge Market Variables}

From Kalshi or prediction markets: probability unemployment $> 5.5\%$, 
recession probability, interest rate changes, CPI surprises. These determine 
hedge payoff structure.

\subsection{Salary Reference Data}

For calibration: BLS wage tables, Glassdoor, Levels.fyi (tech). 
Salary distribution by job level informs income loss magnitude.

\subsection{Unemployment Duration}

From BLS (via FRED): average weeks unemployed (UEMPMEAN). 
Typical mean duration $\approx 20$ weeks. Used for reemployment hazard modeling.

\section{Hazard Model for Job Loss}

We estimate job-loss probability via logistic regression:
\[
P(\text{JobLoss}) = \sigma\left(\beta_0 + \beta_1 U_t + \beta_2 r_t + \beta_3 I_t\right)
\]
where
\[
\sigma(x) = \frac{1}{1+e^{-x}}
\]
and $U_t$ is unemployment, $r_t$ is the interest rate, and $I_t$ is industry risk.
Training uses historical layoffs and macro variables. The weekly Poisson intensity $\lambda$ 
is derived from the monthly probability.

Example coefficient output used in documentation:
\[
\alpha = 0.002,\quad \beta_1 = 0.04,\quad \beta_2 = 0.01,\quad \beta_3 = 0.03
\]

\section{Stochastic Model}

Once $\lambda$ is estimated: $N_t \sim \text{Poisson}(\lambda t)$. Job loss occurs when 
a Poisson jump happens. Employment: $E_t = 1$ if employed, $E_t = 0$ if unemployed.
Income: $dW_t = S E_t \, dt$.

\section{Hedge Payoff Model}

Kalshi contracts are binary: payoff \$1 if event occurs, \$0 otherwise. Example: 
``Unemployment $> 6\%$'' with market price 0.30 implies 30\% probability. 
Hedge value: $H = \text{contract\_price} \times \text{contracts}$.

\section{Monte Carlo Simulation}

We simulate thousands of career paths (default: 5,000 paths, 10-year horizon,
weekly time steps). For each path:

\begin{enumerate}
\item Simulate macro environment (unemployment path)
\item Compute job-loss hazard $\lambda(t)$ from macro state
\item Simulate Poisson event (job loss)
\item Compute salary path $dW_t = S E_t \, dt$
\item Apply hedge payoff
\end{enumerate}

\subsection{Default Dashboard Configuration}

For the default dashboard profile (Tech, Startup, Entry):
\[
S = 120000,\quad N = 5000,\quad T = 10
\]
and Monte Carlo generates
\[
W_t^{(1)}, W_t^{(2)}, \ldots, W_t^{(5000)}.
\]

\section{Risk Metrics}

We evaluate:

\begin{itemize}
\item \textbf{Variance}: income volatility
\item \textbf{Expected Shortfall}: average loss in worst 5\% of cases
\item \textbf{Tail probability}: $P(\text{income drop} > 50\%)$
\end{itemize}

Goal: reduce tail risk while preserving expected income.

\section{Optimal Hedge Ratio}

Choose hedge size using the minimum-variance hedge:
\[
h^* = \frac{\text{Cov}(\text{Loss}, \text{Hedge})}{\text{Var}(\text{Hedge})}
\]
This minimizes the variance of (Loss $- h \cdot$ Hedge).

\section{System Pipeline}

\begin{center}
\begin{tabular}{c}
personal data \\
$\downarrow$ \\
hazard model \\
$\downarrow$ \\
job-loss simulation (Poisson) \\
$\downarrow$ \\
Monte Carlo simulation \\
$\downarrow$ \\
income distribution \\
optimal hedge ratio \\
$\downarrow$ \\
hedge portfolio
\end{tabular}
\end{center}

\section{Discussion}

Prediction markets provide a novel mechanism for hedging personal income risk. 
Unlike traditional insurance, these markets allow individuals to dynamically 
hedge macroeconomic shocks that increase unemployment risk.

The proposed framework combines hazard modeling, stochastic income processes, 
and derivative-style hedging to reduce tail risk in individual wealth 
distributions.

Future work could incorporate industry-specific job-loss models, copula-based 
correlation structures, and real-time hedge optimization using market-implied 
probabilities.

\section{Conclusion}

This paper presents a stochastic framework for hedging job-loss risk using 
prediction market contracts. By modeling employment as a macro-sensitive hazard 
process and constructing correlated hedges, individuals can reduce the downside 
risk of income shocks while maintaining expected earnings.

\end{document}